\section*{अध्याय \ref{ch:b2:atmospheric-thermodynamics} --- वायुमंडलीय ऊष्मागतिकी}

\begin{solution}{exo:b2:atmospheric-thermodynamics:1}
(क) $e = e_s(\qty{278.15}{K}) = \qty{8.73}{hPa}$ और $e_s(\qty{283.15}{K}) = \qty{12.32}{hPa}$।
(ख) $8.73/12.32 = \qty{70.9}{\percent}$। (ग) $r = 0.622 \times 8.73/(850 - 8.73) =
\qty{6.45}{g/kg}$ और $q = r/(1 + r) = \qty{6.41}{g/kg}$।
\end{solution}

\begin{solution}{exo:b2:atmospheric-thermodynamics:2}
(क) $\mathrm{d}\ln e_s/\mathrm{d}T = 5419/T^2$: प्रति केल्विन \qty{7.26}{\percent},
\qty{273.15}{K} पर; \qty{6.31}{\percent}, \qty{293.15}{K} पर; \qty{5.90}{\percent},
\qty{303.15}{K} पर। (ख) $e_s(\qty{20}{\celsius}) = \qty{23.66}{hPa}$ और $e_s(\qty{25}{\celsius})
= \qty{32.26}{hPa}$; अनुपात 1.363, जबकि $1.063^5 = 1.357$। दर
$\ln e_s$ का अवकलज है, इसलिए \qty{5}{K} पर गुणक $\exp(5 \times 0.0620) = 1.363$ है, जहाँ 0.0620
मध्य-बिंदु \qty{22.5}{\celsius} पर की दर है; \qty{6.3}{\percent} को पाँच बार चक्रवृद्धि करना निकट है
पर सटीक नहीं। (ग) $23.66/23.39 = 1.012$: सूत्र \qty{1.2}{\percent} अधिक देता है, क्योंकि
वह त्रिक बिंदु की गुप्त ऊष्मा रखता है, जबकि तापमान बढ़ने पर गुप्त ऊष्मा घटती है; सूत्र का ढाल
बहुत तीखा है, और उसकी त्रुटि
\qty{0}{\celsius} से दूरी के साथ बढ़ती है।
\end{solution}

\begin{solution}{exo:b2:atmospheric-thermodynamics:3}
(क) $125 \times (30 - 14) = \qty{2000}{m}$। (ख) $30 - 9.8 \times 2.0 = \qty{10.4}{\celsius}$।
(ग) $14 - 1.8 \times 2.0 = \qty{10.4}{\celsius}$: तापमान और ओसांक मिलते हैं, जैसा कि संतृप्ति के
तल पर उन्हें मिलना ही चाहिए; नियम यही गणना है। (घ) $125 \times (33 - 14) = \qty{2375}{m}$:
दोपहर जैसे-जैसे भूमि के निकट की वायु को गर्म करती है, मेघ-आधार ऊपर उठते हैं।
\end{solution}

\begin{solution}{exo:b2:atmospheric-thermodynamics:4}
(क) ओसांक वक्र भूमि से लगभग \qty{908}{hPa} तक तापमान वक्र पर पड़ा है:
सबसे निचला आधा किलोमीटर संतृप्त है, या लगभग। (ख) लगभग 908 और \qty{889}{hPa} के बीच
दाब घटने के साथ तापमान वक्र दाईं ओर, गर्म मानों की ओर, मुड़ता है: तापमान
ऊँचाई के साथ बढ़ता है। (ग) LCL लगभग \qty{964}{hPa} पर, कुछ दसियों मीटर ऊपर; LFC लगभग
\qty{730}{hPa} पर; EL लगभग \qty{235}{hPa} पर। (घ) व्युत्क्रमण के शीर्ष पर ओसांक तापमान से
\qty{11}{K} नीचे है, और \qty{20}{K} नीचे \qty{850}{hPa} पर: गर्म, बहुत शुष्क वायु ठंडी, नम वायु के
ऊपर पड़ी है। उसकी सापेक्ष आर्द्रता कम है, इसलिए उसने अपनी गर्माहट संघनन से नहीं पाई; नीचे उतरती
वायु अपना \omterm{def:b2:atmospheric-thermodynamics:moist-air}{मिश्रण-अनुपात} बनाए रखते हुए शुष्क रुद्धोष्म दर से संपीड़न द्वारा गर्म होती है,
जिससे ऐसी परत बनेगी।
\end{solution}

\begin{solution}{exo:b2:atmospheric-thermodynamics:5}
(क) $T_v = 303.15 \times (1 + 0.020/0.622)/1.020 = \qty{306.76}{K}$; सन्निकटन देता है
$303.15 \times (1 + 0.61 \times 0.020) = \qty{306.85}{K}$। (ख) $\rho = p/(R_dT_v) = 10^5/(287.0
\times 306.76) = \qty{1.1358}{kg/m^3}$, जबकि $10^5/(287.0 \times 303.15) =
\qty{1.1494}{kg/m^3}$ शुष्क वायु के लिए: \qty{1.2}{\percent} हल्की। (ग) $T_v - T =
\qty{3.6}{K}$ से। (घ) उठते वायु-पिंड को केवल अपने परिवेश से हल्का होना चाहिए: गर्म समुद्रों के ऊपर
वायु की नमी कई केल्विन के तुल्य उत्प्लावन जोड़ती है, उतना जितना तापमान का
एक उल्लेखनीय आधिक्य।
\end{solution}

\begin{solution}{exo:b2:atmospheric-thermodynamics:6}
(क) $e_s$ = 6.11, 23.66 और \qty{43.54}{hPa}; $r_s = 0.622\,e_s/(1000 - e_s)$ = 3.82, 15.07 और
\qty{28.31}{g/kg}; $\Gamma_s$ = 6.48, 4.20 और \qty{3.44}{K/km}। (ख) शुष्क वायु-पिंड
\qty{9.8}{K} ठंडा होता है, संतृप्त \qty{4.2}{K}: संतृप्त वायु-पिंड \qty{5.6}{K} अधिक गर्म है।
(ग) ऊपर वायु अधिक ठंडी है, $r_s$ छोटा है, संघनित होने को कम वाष्प बचती है, और
$\Gamma_s$ फिर $\Gamma_d$ की ओर लौटता है।
\end{solution}

\begin{solution}{exo:b2:atmospheric-thermodynamics:7}
(क) $e = e_s(\qty{291.65}{K}) = \qty{21.52}{hPa}$, $r = 0.622 \times 21.52/(971.1 - 21.52) =
\qty{14.1}{g/kg}$; $\theta = 292.15 \times (1000/971.1)^{0.2857} = \qty{294.6}{K}$; $\theta_e =
294.6 \exp[2.501 \times 10^6 \times 0.0141/(1004.5 \times 291.5)] = 294.6 \times 1.128 =
\qty{332.3}{K}$। (ख) $e_s(\qty{264.25}{K}) = \qty{3.13}{hPa}$, $r_s = 0.622 \times 3.13/(500 -
3.13) = \qty{3.92}{g/kg}$; $\theta = 264.25 \times 2^{0.2857} = \qty{322.1}{K}$; वायु-पिंड
संतृप्त है, इसलिए $T_{\text{LCL}} = T$ और $\theta_e = 322.1 \exp[2.501 \times 10^6 \times
0.00392/(1004.5 \times 264.25)] = \qty{334.3}{K}$। (ग) $\theta$ \qty{27.5}{K} बढ़ा, संघनित हुई
\qty{10.2}{g/kg} वाष्प की गुप्त ऊष्मा से गर्म होकर; $\theta_e$ \qty{0.6}{\percent} बदला,
अपने सन्निकट सूत्र की परिशुद्धता के भीतर: वह संरक्षित है।
\end{solution}

\begin{solution}{exo:b2:atmospheric-thermodynamics:8}
(क) $\text{CAPE} = 287.0 \times 4 \times \ln(700/300) = 287.0 \times 4 \times 0.847 =
\qty{973}{J/kg}$। (ख) $\sqrt{2 \times 973} = \qty{44.1}{m/s}$। (ग) 1 और
\qty{1500}{J/kg} के बीच: धनात्मक, पर बड़ा नहीं। (घ) दो वक्रों के बीच स्थिर क्षैतिज चौड़ाई की एक
पट्टी; \cref{prop:b2:atmospheric-thermodynamics:areas} के अनुसार उसका क्षेत्रफल $\int
\Delta T\,\mathrm{d}\ln p = 4 \ln(700/300)$ के समानुपाती है, समताप रेखाओं का झुकाव जो भी हो।
\end{solution}

\begin{solution}{exo:b2:atmospheric-thermodynamics:9}
(क) $D_c = \sqrt{3 \times 0.61 \times 80^3/2.1}\,\unit{nm} = \qty{668}{nm} =
\qty{0.668}{\micro m}$; $s_c = 2A/(3D_c) = 2 \times 2.1/(3 \times 668) = 0.00210 =
\qty{0.210}{\percent}$। (ख) $s_c = \sqrt{4A^3/(27\kappa D_d^3)}$ को $D_d$ के लिए हल करने पर: $D_d =
\bigl(4A^3/(27\kappa s_c^2)\bigr)^{1/3}$; \qty{63}{nm}, \qty{0.3}{\percent} पर, और \qty{131}{nm},
\qty{0.1}{\percent} पर। (ग) 63 और \qty{131}{nm} के बीच का हर कण अब सक्रियित होता है: अधिक
बूँदें, और, उतने ही द्रव जल के लिए, छोटी।
\end{solution}

\begin{solution}{exo:b2:atmospheric-thermodynamics:10}
(क) $w = \sqrt{2 \times 130} = \qty{16.1}{m/s}$, भूमि के निकट स्थिर वायु के लिए असंभव।
(ख) व्युत्क्रमण के शीर्ष का विभव तापमान $293.65 \times (1000/889.1)^{0.2857} =
\qty{303.7}{K}$ है; इस $\theta$ वाले सतही वायु-पिंड का $T = 303.7 \times (971.1/1000)^{0.2857}
= \qty{301.2}{K}$, \qty{28.0}{\celsius} है। (ग) भूमि से $125 \times (28.0 - 18.5) = \qty{1190}{m}$
ऊपर, पुराने व्युत्क्रमण-शीर्ष से ऊपर। (घ) एक वाताग्र, जो अपनी ढलान पर वायु को उठाता है; ऊपरी वायु के किसी
विक्षोभ के आगे वायु का धीमा उत्थान, जो स्वयं ढक्कन को ठंडा करता है।
\end{solution}

\begin{solution}{exo:b2:atmospheric-thermodynamics:11}
(क) जब $r_s \to 0$, अंश $g$ की ओर और हर $c_p$ की ओर प्रवृत्त होता है। (ख) लिखिए
$L^2r_s/(R_vT^2) = c_p \times \bigl(Lr_s/(R_dT)\bigr) \times LR_d/(R_vTc_p)$, जहाँ
$LR_d/(R_vTc_p) = 5.1$, \qty{303}{K} पर: जब $Lr_s/(R_dT) \gg 1$, तब $L^2r_s/(R_vT^2) \gg c_p$ भी,
और $\Gamma_s \to g\,(Lr_s/(R_dT))/(L^2r_s/(R_vT^2)) = gR_vT/(LR_d) = 9.81 \times 461.5 \times
303/(2.501 \times 10^6 \times 287.0) = \qty{1.91}{K/km}$। (ग) $T_d = T$ के साथ,
$\Gamma_{\text{ओस}} = gR_vT/(LR_d)$, वही। उस सीमा में गुप्त ऊष्मा इतनी बड़ी है कि
वायु-पिंड अपने संतृप्ति-बिंदु से बहुत ऊपर गर्म हुए बिना अपनी वाष्प का लगभग कुछ भी संघनित नहीं कर
सकता: वह अपना \omterm{def:b2:atmospheric-thermodynamics:moist-air}{मिश्रण-अनुपात} बनाए रखता है और स्थिर \omterm{def:b2:atmospheric-thermodynamics:moist-air}{मिश्रण-अनुपात} की रेखा का अनुसरण करता है, जिसके
अनुदिश तापमान $\Gamma_{\text{ओस}}$ की दर से गिरता है। (घ) \qty{30}{\celsius} और \qty{1000}{hPa} पर,
$Lr_s/(R_dT) = 0.81$, 1 से बहुत अधिक नहीं, और $L^2r_s/(R_vT^2) = 4180$, $c_p$ का चार गुना:
$\Gamma_s = \qty{3.4}{K/km}$, अब भी \qty{1.9}{K/km} से बहुत दूर। सीमा तक नहीं पहुँचा जाता।
\end{solution}

\begin{solution}{exo:b2:atmospheric-thermodynamics:12}
(क) $2.868/2.599 = 1.104$ और $1.257/1.032 = 1.218$: \qty{10.4}{\percent} और
\qty{21.8}{\percent}। (ख) बूँदें \omterm{def:b2:atmospheric-thermodynamics:microphysics}{अतिसंतृप्ति} के एक प्रतिशत के कुछ दसवें भाग पर सक्रियित होती हैं,
दसियों गुना कम। (ग) क्रिस्टल निक्षेपण से तेज़ी से बढ़ते हैं; वे वायु से वाष्प लेते हैं,
जो बूँदों के ऊपर संतृप्ति से नीचे गिर जाती है; बूँदें वाष्पित होती हैं और उनकी वाष्प
हिम पर जाती है (बर्जरॉन प्रक्रम), जब तक मेघ पूरी तरह हिममय न हो जाए। (घ) \qty{0}{\celsius} पर दोनों
संतृप्त \omterm{def:b2:atmospheric-thermodynamics:moist-air}{वाष्प दाब} बराबर होते हैं: वे त्रिक बिंदु पर मिलते हैं, और अनुपात उसके निकट
एक की ओर प्रवृत्त होता है।
\end{solution}

\begin{solution}{pb:b2:atmospheric-thermodynamics:1}
\textbf{1.} $e = e_s(\qty{291.65}{K}) = \qty{21.5}{hPa}$।
\textbf{2.} $r = 0.622 \times 21.5/(971.1 - 21.5) = \qty{14.1}{g/kg}$।
\textbf{3.} $q = 14.1/1.0141 = \qty{13.9}{g/kg}$।
\textbf{4.} $e_s(\qty{292.15}{K}) = \qty{22.2}{hPa}$; $21.5/22.2 = \qty{97}{\percent}$।
\textbf{5.} $0.61 \times 0.0141 \times 292.15 = \qty{2.5}{K}$।
\textbf{6.} $125 \times 0.5 = \qty{63}{m}$; $p = 971.1 \exp[-9.81 \times 62.5/(287.0 \times 294.7)]
= \qty{964}{hPa}$।
\textbf{7.} $\theta = 292.15 \times (1000/971.1)^{0.2857} = \qty{294.6}{K}$।
\textbf{8.} $16.9 - 17.3 = \qty{-0.4}{K}$: अधिक ठंडा।
\textbf{9.} $(20.5 - 16.4)/(1102 - 924) = 4.1/178$: तापमान \qty{23}{K/km} की दर से बढ़ता है,
अर्थात् तापह्रास दर \qty{-23}{K/km} है। एक व्युत्क्रमण: पूर्णतः स्थिर, एक ढक्कन।
\textbf{10.} $15.4 - 20.5 = \qty{-5.1}{K}$।
\textbf{11.} $e_s(\qty{-1.8}{\celsius}) = \qty{5.4}{hPa}$ और $e_s(\qty{18.6}{\celsius}) =
\qty{21.7}{hPa}$: \qty{25}{\percent}। गर्म और शुष्क, यह ऐसी वायु हो सकती है जो नीचे उतरी और संपीड़न
से गर्म हुई, या पवन की आने वाली दिशा में शुष्क भूमि पर गर्म हुई वायु; यह नीचे की नम वायु नहीं है।
\textbf{12.} $18.6 + 18.4 = \qty{37.0}{K}$, $5725 - 1489 = \qty{4236}{m}$ में, \qty{8.73}{K/km}: $\Gamma_s$ और
$\Gamma_d$ के बीच, सप्रतिबंध अस्थिर।
\textbf{13.} $-18.4 - (-8.9) = \qty{-9.5}{K}$: चरम अस्थिरता ($-8$ से $-10$)।
\textbf{14.} भूमि से ऊपर की ओर: 0, $-0.4$, $-0.2$, $-5.1$, $-4.9$, $+1.6$, $+9.5$, $+10.0$,
$+5.6$, $+1.5$, $-9.5$~\unit{K}।
\textbf{15.} 850 ($-4.9$) और \qty{700}{hPa} ($+1.6$) के बीच, भिन्न $4.9/6.5 = 0.754$ पर:
$\ln p = \ln 850 + 0.754 \ln(700/850)$, $p = \qty{734}{hPa}$; $z = 1489 + 0.754 \times 1621 =
\qty{2711}{m}$, भूमि से \qty{2.37}{km} ऊपर।
\textbf{16.} 250 ($+1.5$) और \qty{200}{hPa} ($-9.5$) के बीच, भिन्न $1.5/11.0 =
0.136$ पर: $p = \qty{242.5}{hPa}$ और $z = 10560 + 0.136 \times 1446 = \qty{10757}{m}$, सूचित क्षोभसीमा से
लगभग \qty{200}{m} नीचे (वह \qty{10954}{m} पर है)।
\textbf{17.} खंड (माध्य $\Delta T$ गुणा दाब-अनुपात का $\ln$): $0.80 \times 0.0474
= 0.038$; $5.55 \times 0.3365 = 1.867$; $9.75 \times 0.2231 = 2.176$; $7.80 \times 0.2877 =
2.244$; $3.55 \times 0.1823 = 0.647$; $0.75 \times 0.0305 = 0.023$; योग \qty{6.995}{K}।
$\text{CAPE} = 287.0 \times 6.995 = \qty{2010}{J/kg}$, तीन सार्थक अंकों तक।
\textbf{18.} 1500 और \qty{2500}{J/kg} के बीच: बड़ा।
\textbf{19.} $-0.2 \times 0.0486 - 0.3 \times 0.0189 - 2.65 \times 0.0207 - 5.0 \times 0.0450 -
2.45 \times 0.1467 = -0.655$~\unit{K}; $\text{CIN} = 287.0 \times (-0.655) = \qty{-188}{J/kg}$।
\textbf{20.} तालिका केवल ग्यारह तल रखती है: उनके बीच समलंब नियम वक्रों को
$\ln p$ में सीधे खंडों से बदल देता है और उनके मोड़ चूक जाता है। \qty{1}{hPa} का ग्रिड वक्रों का अनुसरण करता है। \omterm{def:b2:atmospheric-thermodynamics:virtual}{आभासी तापमान}
वायु-पिंड की जलवाष्प का उत्प्लावन जोड़ता है, जो संतृप्त वायु-पिंड में उसके आसपास की
शुष्क वायु से अधिक प्रचुर है: बड़ा CAPE, छोटा CIN।
\textbf{21.} $\sqrt{2 \times 2010} = \qty{63}{m/s}$।
\textbf{22.} $10757 - 2711 = \qty{8046}{m}$, $63.4/2 = \qty{31.7}{m/s}$ पर: \qty{254}{s}, लगभग चार
मिनट।
\textbf{23.} ऊर्ध्व धारा द्वारा ढोए जाने वाले संघनित जल का भार; उसके आसपास की अधिक शुष्क वायु का
भीतर मिलना (अंतर्ग्रहण), जो बूँदों के वाष्पन से उसे ठंडा करता है।
\textbf{24.} वायु-पिंड सिद्धांत इस वायु से पोषित तूफ़ान की ऊर्ध्व धारा को लगभग
\textbf{\qty[text-series-to-math=true]{63}{m.s^{-1}}} की गति देता है, एक ऊपरी सीमा।
\end{solution}
